Le nitrate d'ammonium $\ce{NH4NO3}$ est un composé ionique présent dans divers
engrais. Un sac d'engrais porte l'indication suivante~: <<Pourcentage en masse
75\% de nitrate d'ammonium>> Pour vérifier le pourcentage massique en nitrate
d'ammonium indiqué par le producteur, on prépare une solution aqueuse
$\mathrm{(S_{A})}$ par dissolution de la masse $m=\qty{15.0}{\g}$ d'engrais
dans le volume $V_{0}=\qty{1.0}{\L}$ d'eau distillée.

On dose les ions ammonium $\ce{NH^+_{4(aq)}}$ présent dans un volume
$V_{A}=\qty{10.0}{\mL}$ de la solution $\mathrm{(S_{A})}$ par une solution
aqueuse $\mathrm{(S_{B})}$ d'hydroxyde de sodium
$\ce{Na^+_{(aq)} + HO^-_{(aq)}}$ de concentration molaire
$C_{B}=\qty{0.10}{\mol\per\L}$. Le volume de la solution $\mathrm{(S_{B})}$
versé à l'équivalence est $V_{B,\mathrm{E}}=\qty{14.0}{\mL}$.

\begin{donnees}
    \begin{itemize}
        \item $M(\ce{NH4NO3})=\qty{80.0}{\g\per\mol}$
    \end{itemize}
\end{donnees}

\begin{enumerate}
    \item Écrire l'équation de la réaction qui se produit entre les ions
          ammonium $\ce{NH^+_{4(aq)}}$ et les ions hydroxyde $\ce{HO^-_{(aq)}}$
          au cours du dosage, sachant qu'elle est totale.
    \item Déterminer la valeur de la concentration molaire $C_{A}$ des ions
          ammonium $\ce{NH^+_{4(aq)}}$ dans la solution $\mathrm{(S_{A})}$.
    \item Le pourcentage massique en nitrate d'ammonium contenu dans cet engrais
          s'exprime par la relation~: $\frac{m(\ce{NH4NO3})}{m}$ avec $m$
          la masse de l'engrais.
    
          Calculer le pourcentage massique en masse de nitrate d'ammonium
          contenu dans cet engrais. Comparer à la valeur annoncée par le
          fabriquant.
\end{enumerate}

